\section{$A_0$: Motion and Time as Ontological Horizon}

\subsection{$A_0$ as First Actuality}

The actualization chain begins not with perception but with the ontological conditions perception presupposes. $A_0$ designates the chain's first actuality: motion (\textit{kinēsis}, \gk{κίνησις} `motion, change') and time (\textit{chronos}, \gk{χρόνος} `time'), the co-implicated structures whose prior actuality is the necessary condition for everything that follows. Before sensible objects can exercise their sensible powers, before sense-organs can hold first actuality as capacity, before the rational soul can actualize the counting through which temporal succession becomes articulated, motion and time must already be in act. The chain cannot begin with a static potentiality awaiting external awakening; it must begin with an actuality that is itself a process. The chapter's notation is settled: $A_n$ names actualization stages (nodes), and $M_n \rightarrow M_{n+1}$ names the transitions between them. The chain $A_0 \rightarrow A_1 \rightarrow A_2 \rightarrow A_3 \rightarrow A_4$ runs from this first actuality through perception, \textit{phantasia} (\gk{φαντασία} `imagination'), articulated desire, and the \textit{orektikon} (\gk{ὀρεκτικόν} `desiderative') transition to overt action.

The architectural claim that motion and time are jointly the first actuality rests on Aristotle's doctrine of the priority of actuality over potentiality in three senses: in account (\textit{logos}, \gk{λόγος} `account, reason'), in time, and in substance (\textit{ousia}, \gk{οὐσία} `substance, being'). Aristotle defends this priority in \textit{Metaphysics} IX.8: the potential is intelligible only through the actuality it is potential \textit{for} (1049b12--17); any individual potentiality presupposes a prior actuality capable of bringing it into play (1049b17--1050a3); the end (\textit{telos}, \gk{τέλος} `end, goal') is more fully real than the process directed toward it (1050a4--16). The threefold priority licenses the chain's iterated structure: each completed actuality serves as the unmoved originator of the motion that produces the next, and Aristotle insists, in a passage Burke marks as pivotal, that ``one actuality always precedes another in time right back to the actuality of the eternal prime move'' (Burke, \textit{Grammar}, p.~253, glossing \textit{Metaphysics} IX.8). The iteration is licensed at every node by the perpetual priority of actuality, and Burke isolates the formal dimension: the priority is not first a temporal claim but a formal one, in which ``man is `prior' to boy because man has already attained its complete form whereas boy has not'' (Burke, \textit{Grammar}, pp.~261--262). The formal priority licenses the iterated chain by which each completed actuality functions as both \textit{telos} for what preceded it and \textit{archē} (\gk{ἀρχή} `principle, origin') for what follows. $A_0$ is the limit case: an actuality without prior actuality, since motion and time are co-eternal. Aristotle states the co-eternity explicitly: ``it is impossible that movement should either come into being or cease to be; for it must always have existed'' (\textit{Metaphysics} XII.6, 1071b6--8), and the same passage yokes motion and time: ``nor can time come into being or cease to be; for there could not be a before and an after if time did not exist. Movement also is continuous, then, in the sense in which time is; for time is either the same thing as movement or an attribute of movement'' (1071b8--11).

The chapter's architecture must be read against its twofold project. The complete actualization chain the chapter develops is at once the perception-to-movement chain of \textit{De Motu Animalium} --- sense-perception or thought $\rightarrow$ imagination $\rightarrow$ desire $\rightarrow$ affections $\rightarrow$ organic parts $\rightarrow$ movement (\textit{Movement of Animals} 702a17--19) --- and the actualization chain of desire (\textit{orexis}, \gk{ὄρεξις} `desire') in which the source of all animal locomotion is articulated. These are not two parallel chains. They are one chain. Every cognitive act --- including pure intellection --- is itself a \textit{kinēsis} of the soul, since ``being pained or pleased, or thinking'' are themselves modes of motion of the ensouled being (\textit{De Anima} I.4, 408b5--7); and \textit{orexis} is the moved mover that produces all such motion: ``that which at once moves and is moved'' (\textit{De Anima} III.10, 433b17--18). The chain $A_0 \rightarrow A_4$ thus names one structure: $A_0$ is the motion-time horizon, and the successive stages unfold the single chain through which desire moves the animal via perception, \textit{phantasia}, evaluative articulation, and bodily preparation. The chapter's second purpose --- to demonstrate how Aristotle's \textit{phantasia}, evolved by Heidegger in his analyses of productive imagination and ecstatic temporality, is the human capacity through which time becomes a coherent narrative of past, present, and future (memory, present perception, anticipation) and therefore underwrites every aspect of animal movement --- begins from this $A_0$ horizon, since \textit{phantasia} is the capacity that articulates the motion-time stratum into a coherent narrative manifold.

The Heideggerian register operates from the section's opening, not adjacent to it. Heidegger glosses Aristotelian motion in the \textit{Basic Concepts of Aristotelian Philosophy} as the structural-existential determination of a being's \textit{there}: if a being is determined in its \textit{there} as Aristotle determines \textit{kinēsis}, then it is in movement; \textit{kinēsis} accordingly names what Heidegger develops as \textit{Bewegtheit} (`movedness, being-in-movement'), the way of being of beings whose mode of existence is to be underway toward what they are not yet. The priority of actuality, on this reading, is not a free-standing metaphysical thesis about substance and capacity but a determination of how a comporting being holds itself open toward its ownmost possibilities --- including the possibilities that, at $A_0$, the entire chain of perception-to-action presupposes. The plan of the section follows from this dual register. Motion as \textit{energeia ateles} (\gk{ἐνέργεια ἀτελής} `incomplete actuality, actuality whose end lies beyond it') opens the analysis; the chain structure inherited from priority follows; the relational dyadic structure of \textit{kinēsis} and time as the number of motion develop the analysis through to the soul's constitutive role, at which point the threshold of $A_1$ is reached. The deferrals are flagged at the section's close.

\subsection{\textit{Kinēsis} as \textit{Energeia Ateles}: Being-in-Movement}

Aristotle's compact definition of motion establishes the term whose ontological texture this subsection must specify:
\begin{adjustwidth}{0.5in}{0in}
the fulfilment of what is potentially, as such, is motion --- e.g. the fulfilment of what is alterable, as alterable, is alteration; of what is increasable and its opposite, decreasable (there is no common name for both), increase and decrease; of what can come to be and pass away, coming to be and passing away; of what can be carried along, locomotion. (\textit{Physics} III.1, 201a11--14)
\end{adjustwidth}
Motion is the actuality (\textit{energeia}, \gk{ἐνέργεια} `being-at-work, actuality') of the potential (\textit{dynamis}, \gk{δύναμις} `capacity, potentiality') qua potential; it is a genuine actuality, and yet it is incomplete because its \textit{telos} lies outside itself. Aristotle's pairing in \textit{Metaphysics} IX.6 sharpens the cut: ``at the same time we are seeing and have seen, are understanding and have understood, are thinking and have thought; but it is not true that at the same time we are learning and have learnt, or are being cured and have been cured. \ldots\ For every movement is incomplete --- making thin, learning, walking, building; these are movements, and incomplete movements'' (\textit{Metaphysics} IX.6, 1048b23--30). Seeing and thinking are \textit{energeiai} in the strict sense, since they contain their \textit{telos} within themselves; the moment one has seen, one is seeing, and the activity does not cease on reaching its terminus. Building, learning, and curing are \textit{kinēseis} because their terminus lies beyond the activity itself; the moment the house exists, building has ceased.

The peculiar ontological texture of motion follows from this asymmetry. Motion is the kind of actuality that is at once genuinely present --- the builder is at work, the house is rising --- and incomplete in respect of what it is becoming. Hawhee observes of this register that \textit{energeia} ``names something like physical presence: \textit{energeia} means the presence of the thing'' (Hawhee, ``Rhetorical Vision,'' p.~154). If \textit{energeia} names the presence of the thing, then \textit{kinēsis} as \textit{energeia ateles} names a peculiar mode of presence: the presence of something that is not yet fully what it is becoming. The bronze statue example sharpens the point. The motion of bronze becoming a statue is not the fulfilment of bronze \textit{as bronze} but the fulfilment of bronze \textit{as potentially a statue} (\textit{Physics} III.1, 201a28--33). The qua-structure governs everything: what is actualized in motion is actualized \textit{as} potential, not as achieved completion. Aristotle distributes motion across four categories --- substance, quality, quantity, and place --- and finds in each a pair of contraries between which the moving thing travels (\textit{Physics} V.2, 226a23--30); the world is always already in motion, and there is no moment at which potentiality sits inert, waiting to be actualized from without. $A_0$ is always already actual, and it is actual as ongoing.

The two-modes-of-being-affected distinction Aristotle draws in \textit{De Anima} II.5 reinforces the point and prepares the relational analysis to come. ``The expression `to be acted upon' has more than one meaning; it may mean either the extinction of one of two contraries by the other, or the maintenance of what is potential by the agency of what is actual and already like what is acted upon, as actual to potential'' (\textit{De Anima} II.5, 417b2--5). In one sense being-affected is a destruction; in the other it is a preservation in which a potential is brought into its actuality. Hearing a sound is a being-affected of the second sort: my hearing is brought to its full being, not destroyed. The same structural point governs $A_0$: the motion-time horizon is not the kind of actuality that is exhausted in producing the next stage of the chain but the kind that preserves itself as the ground on which the next stage stands.

Heidegger's reading of Aristotelian \textit{kinēsis} in the \textit{Basic Concepts of Aristotelian Philosophy} (the 1924 Marburg lectures, \textit{GA} 18) and in the \textit{Phenomenological Interpretations of Aristotle} fragments (\textit{GA} 61) transposes this analysis from the level of natural-philosophical category to the level of fundamental ontology. Motion, on Heidegger's reading, is a \textit{Seinscharakter} (`character of being'): if a being is determined in its \textit{there} as Aristotle determines \textit{kinēsis}, then that being is in movement, and \textit{kinēsis} accordingly names what Heidegger develops as \textit{Bewegtheit} --- ``\textbf{******}'' (\textit{Basic Concepts of Aristotelian Philosophy}, p.~\textbf{******}).\footnote{Verbatim text and page reference to be supplied manually. The passage appears in \textit{GA} 18 / \textit{Basic Concepts of Aristotelian Philosophy} in Heidegger's articulation of the Aristotelian definition of \textit{kinēsis} as the structural determination of the being-there of Dasein-style beings; Heidegger's term \textit{Bewegtheit} does the work of recovering Aristotelian \textit{energeia ateles} at the level of fundamental ontology. The pertinent material is in the chapter-area on \textit{kinēsis} and being-in-movement, alongside Heidegger's reading of \textit{Metaphysics} $\Theta$ on \textit{dynamis} and \textit{energeia}.} The structural payoff is precise. Dasein --- the being whose way of being is to comport itself toward its own being --- is not a static substance which subsequently enters into changes; it is a being whose mode of existence is being-underway-toward, the ahead-of-itself structure that \textit{Being and Time} \S 41 develops under the name \textit{Sorge} (`care'). ``Dasein is its possibility,'' Heidegger writes (SZ \S 9, H.42; Eng.~67), and the formula recovers Aristotelian motion at the level of fundamental ontology --- not as an external addition to a prior substance but as the way in which the comporting being is what it is. The intermediate-register character of \textit{energeia ateles} --- genuinely actual yet structurally incomplete --- is therefore not a marginal metaphysical curiosity but the primitive ontological pattern from which the entire chain inherits its iterated structure: each completed actuality at $A_1$, $A_2$, $A_3$ is both the terminus of its antecedent motion and the unmoved originator of its successor, an \textit{entelecheia} (\gk{ἐντελέχεια} `having-one's-end-within-oneself') in relation to what precedes it and an \textit{archē} in relation to what follows. The chain is iterated because each of its nodes has the same dual structure $A_0$ establishes as primitive. The temporality of taking care that \textit{Being and Time} \S 68 develops --- the unity of awaiting, making-present, and retaining as the temporal articulation of Dasein's circumspective concern (SZ \S 68b, H.342) --- recovers, at the existential level, what the chain establishes at the level of categorial structure: the comporting being is constituted by the ongoing motion of its being-underway, and each completed phase of its concern is at once the terminus of the motion that produced it and the principle of the motion that follows.

\subsection{Priority of Actuality and the Chain Structure}

The chain structure of $A_0 \rightarrow A_4$ inherits its logic from the priority of actuality over potentiality. Aristotle states the priority in its most concentrated form:
\begin{adjustwidth}{0.5in}{0in}
everything that comes to be moves towards a principle, i.e. an end. For that for the sake of which a thing is, is its principle, and the becoming is for the sake of the end; and the actuality is the end, and it is for the sake of this that the potentiality is acquired. For animals do not see in order that they may have sight, but they have sight that they may see. (\textit{Metaphysics} IX.8, 1050a7--11)
\end{adjustwidth}
The end is not merely the terminus of a process but its \textit{archē}: the completed form furnishes the reason why the process occurs at all and the measure by which its stages are ordered. Sight is not a capacity that becomes actualized as seeing; seeing is the actuality for the sake of which sight exists as a capacity. The three senses of priority each carry architectural consequences for the chain. In account, the potential is intelligible only through the actuality it is potential for; the chain's apparatus is therefore intelligible only through the completed $A_4$ that gives it its shape. In time, any individual potentiality presupposes a prior actuality capable of bringing it into play; $A_0$'s motion-time horizon is what brings into play the natural beings of $A_1$. In substance, the \textit{telos} is more fully real than the process directed toward it; the completed action at $A_4$ is more fully real than the chain's earlier phases, even though those phases must precede it in time. Burke isolates the architectural payoff: the entelechy ``allowed [Aristotle] to introduce another kind of priority, namely the `principle' involved in a given form'' (Burke, \textit{Grammar}, pp.~261--262). The priority licenses an iterated structure in which each completed actuality functions as both terminus and origin.

The principle by which motions are named in the chain is articulated cleanly in the \textit{Physics}: ``it is that to which rather than that from which the motion proceeds that gives its name to the change'' (\textit{Physics} V.1, 224b7--8). Each $M_n \rightarrow M_{n+1}$ in the chain takes its name from its terminus, not from its starting point. The motion of perception is named for the perceptual \textit{energeia} it produces; the motion of \textit{phantasia}-formation, for the \textit{phantasma} (\gk{φάντασμα} `image, appearance') it leaves behind; the \textit{orektikon} transition, for the desire's discharge into action. Aristotle's geometer illustration sharpens the relation between potential and actual:
\begin{adjustwidth}{0.5in}{0in}
the potentially existing constructions are discovered by being brought to actuality (the reason being that thinking is the actuality of thought); so that the potentiality is discovered from actuality (and therefore it is by an act of construction that people acquire the knowledge). (\textit{Metaphysics} IX.9, 1051a29--33)
\end{adjustwidth}
All the possible geometric configurations exist potentially in the figure, but they are not discovered by introspection; the geometer who proves the theorem actualizes the potentialities that were already latent there, and the actualization makes them manifest. The principle generalizes: each stage's potentialities are not abstract capacities awaiting filling-in but structural possibilities that become intelligible only when the prior actuality has licensed them.

The transitional character of $A_0$ as motion-medium is articulated in Aristotle's general formulation of motion as passage between states: ``if there is movement and something moved, that which is moved must first be in that out of which it is to be moved, and then not be in it, and move into the other and come to be in it'' (\textit{Metaphysics} XI.6, 1063a19--21). Motion is the medium of passage; the chain's iterated structure is a sequence of such passages, each instantiating the same transitional ontology. The relation between the chain's stages is therefore not a causal-mechanical sequence but a hermeneutical one. The completed actuality at each stage --- the for-the-sake-of-which --- is always already operative in the understanding that guides the actualization. Aristotle articulates this in his own register as the operation of the \textit{telos}; Heidegger generalizes the structure at the level of Dasein's interpretive comportment as the hermeneutical circle. Understanding, on the analysis of \textit{Being and Time} \S 32, projects itself upon possibilities that are themselves disclosed by the fore-structure of understanding: \textit{Vorhabe} (`fore-having, what one has in advance'), \textit{Vorsicht} (`fore-sight, what one sees in advance'), and \textit{Vorgriff} (`fore-conception, what one grasps in advance') jointly articulate the as-structure of interpretation, in which what is interpreted is already disclosed as something prior to any propositional articulation (SZ \S 32--33, H.150--159). The circle is the structural condition for any actualization that proceeds toward a \textit{telos}: the \textit{telos} must already be operative in the comportment that aims at it, and the comportment must already be guided by the \textit{telos} it does not yet possess. ``What is decisive is not to get out of the circle but to come into it in the right way,'' Heidegger writes (SZ \S 32, H.153; Eng.~195). The chain's iterated structure --- \textit{energeia ateles} generating \textit{entelecheia} generating new conditions for \textit{energeia ateles} --- is the proto-form of the hermeneutical circle Heidegger generalizes for Dasein's interpretive comportment as a whole. Aristotle's apparatus articulates the structural-categorial features of the chain; Heidegger's apparatus articulates how those features are determinations of a comporting being's way of being. The chapter deploys the two registers together.

\subsection{The Relational Structure of \textit{Kinēsis}}

A feature of $A_0$ that the chain inherits at every subsequent stage is the relational character of \textit{kinēsis}. Motion is not a property inhering in a single substance but a relation between mover and moved, between the active and passive powers whose joint actualization constitutes the process of change. Aristotle states the relation directly:
\begin{adjustwidth}{0.5in}{0in}
motion is in the movable. It is the fulfilment of this potentiality by the action of that which has the power of causing motion; and the actuality of that which has the power of causing motion is not other than the actuality of the movable; for it must be the fulfilment of \textit{both}. A thing is capable of causing motion because it \textit{can} do this, it is a mover because it actually \textit{does} it. But it is on the movable that it is capable of acting. Hence there is a single actuality of both alike. (\textit{Physics} III.2, 202a13--20)
\end{adjustwidth}
Motion is irreducibly relational, a joint actualization of two capacities --- the capacity to act and the capacity to be acted upon --- and the \textit{energeia} shared between them is not the consequence of the motion but the motion itself. Aristotle specifies the dyadic structure in categorial terms: ``it is necessary that there should be an actuality of the agent and of the patient. The one is agency and the other patiency; and the outcome and end of the one is an action, that of the other a passion'' (\textit{Physics} III.3, 202a21--24). The agent's actuality is \textit{poiēsis} (\gk{ποίησις} `making, producing'); the patient's actuality is \textit{pathēsis} (\gk{πάθησις} `being-affected, undergoing'). Aristotle clarifies the point with the Thebes-Athens analogy: ``Thebes and Athens are different cities, yet the road from Thebes to Athens and the road from Athens to Thebes are the same'' (\textit{Physics} III.3, 202b13--14). One road, two descriptions. One motion, two formal aspects. The distinction is in account, not in substrate.

The dyadic architecture supplied by $A_0$ propagates through every subsequent stage of the chain. At $A_1$ the sensible object's active potency meets the sense-organ's receptive potency in a single perceptual actualization: perception is itself a kind of alteration, in which the sense-organ is moved by the sensible object and the perceptual act is the joint actualization of the organ's capacity to receive form and the object's capacity to act upon it (\textit{De Anima} II.5, 416b33--417a20; cf. III.2, 425b26--27 on the single-event-under-two-\textit{logoi} structure). The full analysis of this propagation belongs to the section on \textit{aisthēsis} (\gk{αἴσθησις} `perception'), but the structural inheritance is fixed at $A_0$. The same schema iterates at the cognitive-motor stages of the chain: what the realizable good moves at $M_3 \rightarrow A_4$ is \textit{orexis}, the moved mover that is at once agent (of the animal's locomotion) and patient (of the apparent good); what \textit{orexis} in turn moves is the animal via its bodily instrument. The chain is dyadic at every node because $A_0$ establishes dyadic relationality as the primitive structural feature of \textit{kinēsis}. Burke captures the consequence in his grammar of motives: the motion of $A$ by $B$ determines both $A$ and $B$, assigning each a role (patient and agent) and a trajectory (from potentiality to actuality), since ``conditions are likewise contextual, as with the conditions of an organism's existence'' (Burke, \textit{Grammar}, pp.~214--215). The contextual character of determination is precisely what the chain's iterated dyadic structure articulates at every node.

Heidegger's existential radicalization of the relational structure is the analysis of Dasein as constituted by its \textit{In-der-Welt-sein} (`being-in-the-world'). Dasein does not first exist as substance and then enter into relations; it is constituted by relational being-in-the-world, in which entities are encountered as ready-to-hand (\textit{Zuhandenheit} `readiness-to-hand'), as members of equipmental contexts (\textit{Zeug} `equipment') articulated by chains of involvement (\textit{Bewandtnis} `involvement, the with-which-it-has-its-business'). Heidegger states the structural claim in \textit{Being and Time} \S 18: ``\textbf{******}'' (SZ \S 18, H.84; Eng.~\textbf{******}).\footnote{Verbatim text and English pagination to be supplied manually. The passage appears in \textit{Being and Time} \S 18 (``Involvement and Significance: the Worldhood of the World''), in the opening paragraphs articulating how involvement (\textit{Bewandtnis}) is the structural character of entities encountered within the world, and how the totality of involvements (\textit{Bewandtnisganzheit}) is grounded in the for-the-sake-of-which (\textit{Worumwillen}) of Dasein's own existence. The pertinent passage develops the claim that the worldhood of the world is constituted by the involvement-relations rather than by an aggregate of independent substances.} The Aristotelian relational \textit{kinēsis} is, on this reading, the proto-form of Heideggerian worldly involvement: motion is always a being-toward, an articulation between an active power and a passive power that are not first independent substances and then jointly active but are co-constituted in the actuality of the motion they share. The chain's iterated dyadic structure is, at the existential level, the structural-categorial trace of the way a being whose mode of existence is being-in-the-world is constituted by its relations to what it encounters.

\subsection{Time (\textit{Chronos}) as the Number of Motion}

Aristotle's celebrated definition of time establishes the second moment of the $A_0$ horizon:
\begin{adjustwidth}{0.5in}{0in}
time is the number of motion in respect of the before and after (\gk{ἀριθμὸς κινήσεως κατὰ τὸ πρότερον καὶ ὕστερον}). (\textit{Physics} IV.11, 219b1--2)
\end{adjustwidth}
Time is not an independent substance --- not a container within which events transpire, not an autonomous metaphysical entity --- but an aspect of motion: motion insofar as it admits of enumeration. Where there is no motion there is no time, and where motion can be numbered there is time as that numbering. The continuity and directional structure of time are inherited from motion, and motion inherits these from spatial magnitude: ``the movement goes with the magnitude, and time, as we maintain, goes with motion'' (\textit{Physics} IV.11, 219a10--13). Time is accordingly ``either the same thing as movement or an attribute of movement'' (\textit{Metaphysics} XII.6, 1071b10--11).

The definition turns on a distinction that is easy to miss but architecturally load-bearing. ``Time, then, is a kind of number. Number, we must note, is used in two ways --- both of what is counted or countable and also of that with which we count. Time, then, is what is counted, not that with which we count: these are different kinds of things'' (\textit{Physics} IV.11, 219b5--10). Time belongs to the first category: it is the counted number of motion, not the counting activity. The distinction will prove decisive when the soul's role enters the analysis, for it allows time to be located in the structure of motion itself --- as what \textit{can} be counted --- while still requiring a counter for the full actualization of time as numbered. The grounding of the definition is at once structural and experiential. ``We perceive movement and time together,'' Aristotle observes, ``even when it is dark and we are not being affected through the body, if any movement takes place in the mind we at once suppose that some time has indeed elapsed'' (\textit{Physics} IV.11, 219a4--8). Neither motion nor time is epistemically prior to the other; they are apprehended together, and the reciprocity is ontological as well as epistemic.

The now (\gk{τὸ νῦν} `the now') is the principle by which time's continuum is at once articulated and divided. ``The `now' determines time, insofar as time involves the before and after'' (\textit{Physics} IV.11, 219b12--13); ``time, then, also is both made continuous by the `now' and divided at it'' (\textit{Physics} IV.12, 220a5). The now is at once the principle of continuity --- the moving boundary that links successive durations --- and the principle of division --- the point at which past ends and future begins. ``What is \textit{before} is in time; for we say `before' and `after' with reference to the distance from the `now', and the `now' is the boundary of the past and the future'' (\textit{Physics} IV.14, 223a4--6). Aristotle illuminates the now's mode of being through the body-carried-along analogy: ``the `now' corresponds to the body that is carried along, as time corresponds to the motion. For it is by means of the body that is carried along that we become aware of the before and after in the motion, and if we regard these as countable we get the `now''' (\textit{Physics} IV.11, 219b22--25). The moving body is always the same body even as it traverses different positions; the now is always the same now even as it marks different moments. One in substrate, different in being.

Heidegger reads the Aristotelian analysis in \textit{Being and Time}, Division Two, Chapter VI, as the most rigorous articulation of what he calls \textit{Innerzeitigkeit} (`within-time-ness, innertimeness') --- time as publicly encountered, the time-in-which entities and events occur. Aristotle's analysis is correct as far as it goes: it captures the structure of public, world-time, the time within which the chain's actualizations unfold and which is articulated by the now's bounding and joining function. What the analysis does not thematize is the more primordial structure of temporality (\textit{Zeitlichkeit} `temporality') and historicity (\textit{Geschichtlichkeit} `historicality') from which \textit{Innerzeitigkeit} is derived. The derivative time of the now-sequence presupposes Dasein's ecstatic temporalizing as the originary source from which the public-time concept is generated through a levelling-off. Heidegger states the relation in his explicit engagement with Aristotle's definition: ``\textbf{******}'' (SZ \S 81, H.421--422; Eng.~\textbf{******}).\footnote{Verbatim text and English pagination to be supplied manually. The passage appears in \textit{Being and Time} \S 81 (``Within-time-ness and the genesis of the ordinary conception of time''), in Heidegger's explicit engagement with Aristotle's definition of time as the number of motion. Heidegger here characterizes the Aristotelian definition as the most rigorous philosophical account of the time encountered in everyday concern, while indicating that the analysis leaves unthematized the originary temporality from which \textit{Innerzeitigkeit} derives. The pertinent passage is on H.421--422, in the early paragraphs of \S 81 just after the transition from \S 80 on the genesis of the ordinary conception of time, and in proximity to Heidegger's claim that Aristotle's definition fixes the natural philosophy of time at the level of within-time-ness without proceeding to its existential ground.} The structural-philosophical point is precise. Aristotle has captured the \textit{Innerzeitigkeit}-stratum of time without thematizing its grounding in originary temporality; the chapter's apparatus, accordingly, deploys what Heidegger calls \textit{Innerzeitigkeit} as the time-stratum within which the chain unfolds, while the more primordial temporality that the chain-actualizations presuppose is held open by the Heideggerian framework. The two registers operate together. Aristotle's analysis is the most precise account available of the time-in-which the actualization chain proceeds; Heidegger's analysis is the most precise account available of the temporality from which that public time derives.

The implication for the chapter's broader project is direct. The cognition of time as a coherent narrative of past, present, and future --- the temporal manifold within which memory, present perception, and anticipation are unified --- is not delivered by the Aristotelian definition alone. The definition captures countable succession; it does not, on its own, articulate the synthetic activity through which the three temporal modes are knit into a unified phenomenal manifold. That synthetic activity is the work of \textit{phantasia}, which the subsequent sections develop in detail. The Aristotelian motion-time stratum at $A_0$ is the substrate within which \textit{phantasia}'s temporal-synthesizing activity unfolds. Aristotle himself prepares the thesis: ``imagination is a feeble sort of sensation, and memory and expectation involve imagination; therefore things remembered and expected can be pleasant'' (\textit{Rhetoric} I.11, 1370a28--33). \textit{Phantasia} is what allows the soul to be moved by absences, possibilities, appearances, and counterfactuals; it is the capacity that articulates the motion-time substrate into the temporal manifold within which past, present, and future co-temporalize for the comporting being.

\subsection{The Soul's Constitutive Role}

Aristotle's analysis of the soul in \textit{De Anima} establishes the soul's constitutive role in the actualization chain. The soul is ``a substance (\textit{ousia}) in the sense of the form (\textit{eidos}, \gk{εἶδος} `form, kind') of a natural body'' (\textit{De Anima} II.1, 412a19--20); it is the first actuality of a natural body that has life potentially, what makes a body into a living organism capable of perceiving, moving, and thinking. The soul does not exist apart from the body; it is the body's organizing principle, its way of being what it is. Hence the soul is itself an instance of the priority of actuality: it is the form that must be actual before the body's capacities can be exercised. The point connects directly to the motion-time parallel. Just as motion is the actuality of the potential qua potential, so the soul is the actuality of the body's potential for life; just as time is the numbered aspect of motion, so the rational soul is the numbering principle that articulates temporal flow.

Aristotle's three-grade schema of potentiality clarifies the architectural placement of $A_1$: first potentiality is the bare capacity, second potentiality or first actuality is the developed capacity, second actuality is the exercise of that capacity (\textit{De Anima} II.5, 417a21--b2). $A_1$ corresponds to first actuality: the perceiving organism equipped with developed capacities but not yet engaged with any particular sensible object. The hierarchical ordering of psychic capacities (\textit{De Anima} II.3) reinforces the iterated structure: each level presupposes the actualization of those below it --- perception presupposes the nutritive capacity, \textit{phantasia} presupposes perception, thought presupposes \textit{phantasia}.

The architectural weight of $A_0$ becomes most explicit at the soul's relation to time. ``If nothing but soul, or in soul reason, is qualified to count, it is impossible for there to be time unless there is soul, but only that of which time is an attribute --- i.e., if movement can exist without soul'' (\textit{Physics} IV.14, 223a25--27). The conditional is precise. Motion may exist without soul: natural bodies move whether or not anyone observes them. But time --- the numbered aspect of motion --- requires a counter, and where there is no counter there can be no number counted. The chapter adopts a stronger interpretive reading of the conditional than Aristotle's text strictly compels: time without a counting soul is not merely unnumbered but ontologically incomplete --- present as countability, structurally real as the before-and-after of motion, yet not fully actualized as time. The reading is interpretive rather than textually compelled, and the chapter flags it as such. It coheres with Aristotle's broader priority-of-actuality commitments: if time is fully actual only as numbered, what makes the numbering possible is constitutive of time's full actuality. The textual warrant is cautious-conditional, and extended engagement with the alternative reading is deferred.

Heidegger radicalizes the soul-as-counter thesis as originary temporality (\textit{Zeitlichkeit}). Originary temporality is not a sequence of nows but the unity of three ecstases (\textit{Ekstasen} `ecstases, standings-outside'): having-been (\textit{Gewesenheit} `having-been-ness'), present (\textit{Gegenwart} `present, making-present'), and projected-futural (\textit{Zukunft} `future, coming-toward'). Each ecstasis is a mode in which Dasein temporalizes itself; the three are co-original, and their unity is the originary phenomenon from which all derivative time-conceptions are drawn. Heidegger states the relation: ``\textbf{******}'' (SZ \S 65, H.328--329; Eng.~\textbf{******}).\footnote{Verbatim text and English pagination to be supplied manually. The passage appears in \textit{Being and Time} \S 65 (``Temporality as the ontological meaning of care''), in Heidegger's articulation of temporality as the unity of three ecstases (\textit{Gewesenheit}, \textit{Gegenwart}, \textit{Zukunft}) and as the originary phenomenon from which \textit{Innerzeitigkeit} is derived. The pertinent material is in the middle paragraphs of \S 65 around H.328--329, where Heidegger first introduces the three ecstases by name and characterizes their co-original unity, and from which his subsequent argument in \S\S 66--68 develops temporality as the temporalizing-source for the existential structures of being-in-the-world. Cross-reference also \S 74, H.385--386, where Heidegger develops historicity (\textit{Geschichtlichkeit}) as grounded in originary temporality.} The Aristotelian insight that time requires soul as counter is secured at the existential-ontological level by this analysis: Dasein's ecstatic temporality \textit{is} the originary phenomenon from which the public-time of the now-sequence derives. The chapter's interpretive reading of \textit{Physics} IV.14 --- that time without soul is ontologically incomplete --- aligns most cleanly with Heidegger's recovery of Aristotelian time at the existential level, since on the Heideggerian reading the public time of the chain is grounded in Dasein's temporalizing, and a motion-time stratum without any soul to count it would lack the originary temporality that makes it intelligible as time at all.

\subsection{The Threshold of $A_1$}

The threshold of $A_1$ is the architectural pivot at which the chapter's central capacity-thesis enters the analysis. Aristotle's \textit{phantasia} and Heidegger's productive imagination name the same temporal-synthetic capacity through which a finite, embodied, comporting being cognizes time as a coherent narrative of past, present, and future. Aristotle states the foundation: ``all memory implies time; therefore only animals that perceive time can remember, and the organ by which they perceive time is that by which they remember'' (\textit{De Memoria} 449b28--30); and ``imagination is a feeble sort of sensation, and memory and expectation involve imagination'' (\textit{Rhetoric} I.11, 1370a28--30). \textit{Phantasia} is the capacity that ranges across all three temporal modes --- it presents the present object, retains the absent past object as memory, and projects the not-yet object as anticipation. Without \textit{phantasia}, the soul could be moved only by what is perceptually present; with \textit{phantasia}, the soul can be moved by absences, possibilities, appearances, and counterfactuals, which is to say by the full range of objects required for the actualization chain to operate across memory and anticipation rather than within the immediate now alone. Heidegger's analysis of productive imagination in \textit{Kant and the Problem of Metaphysics} (\textit{GA} 3) radicalizes this Aristotelian capacity as the temporalizing-synthesis that knits the three ecstases into the unity of originary temporality. The full development of the \textit{phantasia}-thesis belongs to the subsequent sections; what $A_0$ establishes is the ontological horizon within which \textit{phantasia}'s temporal-synthesizing activity becomes intelligible, and it is by virtue of this synthesizing activity that the chain spans past, present, and future as a coherent narrative manifold.

With the phantasia-thesis named, the threshold of $A_1$ can be stated with precision. $A_0$'s motion-time horizon licenses $A_1$'s natural beings, sensible objects exercising their sensible powers, sense-organs holding their first actuality as capacities, and rational souls capable of articulating temporal succession. Motion being actual is what allows sensible objects to actually exercise their sensible powers --- the color visible, the sound audible, the warm warming. Motion being actual is what allows the medium to transmit the sensible form from object to organ. Time being actual is what allows the temporal continuum within which the transmission unfolds; and time becoming fully actual with the arrival of rational soul is what allows the sequential articulation of perceptual, imaginative, and cognitive acts the subsequent sections will trace. $A_1$ names the world in which all four conditions obtain: rational souls, sensible objects actually exercising their sensible powers, sense-organs holding first actuality as capacity, and a temporal-narrative manifold within which the actualization chain proceeds. Each of these depends on $A_0$, and each will itself become the unmoved originator of further motions as the chain unfolds. The chain begins.

\subsection*{Deferrals and Forward Orientations}

Three matters bear future development and are flagged here for the reader's orientation. First, the interpretive reading of \textit{Physics} IV.14, 223a25--27 adopted in this section --- that time without soul is ontologically incomplete rather than merely unnumbered --- is defensible on priority-of-actuality grounds, but the textual warrant is cautious-conditional. Alternative readings that treat time as fully actual independently of soul merit extended engagement, with Coope's work on Aristotelian time and Broadie's commentary on \textit{Physics} IV as the principal interlocutors. Second, the relation between time as a common sensible (\textit{De Anima} III.1, 425a16), the discriminating faculty's unity across the now (\textit{De Anima} III.2, 426b17--427a1), and the role of \textit{phantasia} in producing continuity across discrete perceptual instants bears on how $A_0$'s time becomes $A_1$'s cognitively articulated temporal manifold. The \textit{De Memoria} passages on memory requiring the recognition of past-qua-past (449b22--25, 450a11--13), together with \textit{De Anima} III.10, 433b5--10 on the sense of time co-implicated with deliberative \textit{phantasia}, specify the psychological apparatus through which the cognitive articulation of time occurs; the systematic treatment of \textit{phantasia} as the temporal-narrative-synthesizing capacity is the chapter's structural pivot and belongs to the sections on \textit{aisthēsis}, \textit{phantasia}, and the orientational modes at $M_2 \rightarrow A_3$. Third, the entelechial structure displayed by the chain at each completed actuality --- the way each $A_n$ functions as both \textit{telos} and \textit{archē} --- suggests an affinity with the rhetorical theory of motivated action subsequent chapters will develop; the full articulation of this affinity lies beyond the present section. With $A_0$ specified, the analysis can now turn to $A_1$, where the natural beings, sensible objects, and rational souls whose first actualities $A_0$'s motion-time horizon licenses begin to enter the chain in their own structural roles.
